\documentclass[a4paper,10pt]{article}
\usepackage[utf8]{inputenc}
\usepackage{amssymb}
\usepackage{amsmath}
\usepackage{graphicx}
\newcommand{\HRule}{\rule{\linewidth}{0.5mm}}
\usepackage{subfigure}
\usepackage{multicol}
\usepackage[usenames,dvipsnames]{color}
\definecolor{darkgray}{rgb}{0.95,0.95,0.95}
\usepackage{listings}
\usepackage{color}
\usepackage{textcomp}
\definecolor{listinggray}{gray}{0.9}
\definecolor{lbcolor}{rgb}{0.9,0.9,0.9}
\lstset{
    backgroundcolor=\color{lbcolor},
    tabsize=4,
    rulecolor=,
    language=C++,
    basicstyle=\scriptsize,
    upquote=true,
    aboveskip={1.5\baselineskip},
    columns=fixed,
    showstringspaces=false,
    extendedchars=true,
    breaklines=true,
    prebreak = \raisebox{0ex}[0ex][0ex]{\ensuremath{\hookleftarrow}},
    frame=single,
    showtabs=false,
    showspaces=false,
    showstringspaces=false,
    identifierstyle=\ttfamily,
    keywordstyle=\color[rgb]{0,0,1},
    commentstyle=\color[rgb]{0.133,0.545,0.133},
    stringstyle=\color[rgb]{0.627,0.126,0.941},
}
\lstset{language=C++}
\lstset{backgroundcolor=\color{darkgray}}

\usepackage[left=1.0cm, right=1.0cm, top=2cm, bottom=4cm]{geometry}
\usepackage{fancyhdr}
\pagestyle{fancy}
\usepackage{lastpage}
\renewcommand\headrulewidth{1pt}
\fancyhead[L]{\textsc{Nachos Étape 2 : Entrées/Sorties}}
\fancyhead[R]{\textsc{Polytech' Grenoble}}
\renewcommand\footrulewidth{1pt}
\fancyfoot[R]{ \textsc{RICM 4}}

\begin{document}
\begin{titlepage}

\begin{center}


% Upper part of the page
\includegraphics[width=0.25\textwidth]{./images/logo}\\[1cm]

\textsc{\LARGE Polytech' Grenoble}\\[1.5cm]

\textsc{\Large RICM 4\`eme ann\'ee}\\[1.2cm]


% Title
\HRule \\[0.4cm]
{ \huge \bfseries NachOS\\[0.6cm]
Etape 2: Entrées/Sorties}
\\[0.4cm]

\HRule \\[2cm]

% Author and supervisor
\begin{minipage}{0.4\textwidth}
\begin{flushleft} \large
\emph{\'Etudiants:}\\
Elizabeth \textsc{Paz} \\
Salem \textsc{Harrache}
\end{flushleft}
\end{minipage}
\begin{minipage}{0.4\textwidth}
\begin{flushright} \large
\emph{Professeur:} \\
Vania \textsc{Marangozova}
\end{flushright}
\end{minipage}

\vfill

% Bottom of the page
{\large  février 2012}

\end{center}

\end{titlepage}

\section{Entrées/Sorties asynchrones}

l'objet Console est asynchrones, dans le programme de test (ConsoleTest) on
gère manuellement la syncronisation avec deux sémphores (un pour l'écriture et
un autre pour la lecture) et des handlers qui seront appellés par le traitant
une fois la tâche de lecture/ecriture effectuée. Dans cette partie, on ajoute la
prise en compte du caractère de fin de fichier (EOF).

\textit{code/usrprog/progtest.cc}
\begin{lstlisting}
void ConsoleTest (char *in, char *out)
{
    char ch;

    console = new Console (in, out, ReadAvail, WriteDone, 0);
    readAvail = new Semaphore ("read avail", 0);
    writeDone = new Semaphore ("write done", 0);

    for (;;) {
        readAvail->P (); // wait for character to arrive
        ch = console->GetChar ();
        if (ch == EOF or ch =='q')
            return;
        console->PutChar('<');
        writeDone->P ();
        console->PutChar(ch); // We echo it
        writeDone->P ();
        console->PutChar('>');
        writeDone->P (); // wait for write to finish
    }
}
\end{lstlisting}

Comme on le voir, a chaque fois qu'on écrit ou lit un caractère il faut
utiliser le sémaphore explicitement pour obliger le programme à se bloquer et
reprendre qu'une fois le caractère lu/écrit

\section{Entrées/Sorties synchrones}
On va maintenant gérer la syncronisation directement dans la Console. Pour ça
on va implementer la classe SynchConsole qui aura utilisé de façon
transparentlese les sémphores. On ajoute les deux fichiers
\textit{code/userprog/synchconsole.h} et \textit{code/userprog/synchconsole.cc}

\textit{code/userprog/synchconsole.h}
\begin{lstlisting}
#ifndef SYNCHCONSOLE_H
#define SYNCHCONSOLE_H

#include "copyright.h"
#include "utility.h"
#include "console.h"

class SynchConsole {
    public:
        SynchConsole(char *readFile, char *writeFile);
        // initialize the hardware console device
        ~SynchConsole();
        // clean up console emulation
        void SynchPutChar(const char ch); // Unix putchar(3S)
        char SynchGetChar(); // Unix getchar(3S)
        void SynchPutString(const char *s); // Unix puts(3S)
        void SynchGetString(char *s, int n); // Unix fgets(3S)
    private:
        Console *console;
};
#endif // SYNCHCONSOLE_H
\end{lstlisting}

\textit{code/userprog/synchconsole.cc}
\begin{lstlisting}
#include "copyright.h"
#include "system.h"
#include "synchconsole.h"
#include "synch.h"


static Semaphore *readAvail;
static Semaphore *writeDone;

static void ReadAvail(int arg) {
    readAvail->V();
}

static void WriteDone(int arg) {
    writeDone->V();
}

SynchConsole::SynchConsole(char *readFile, char *writeFile) {
    /* Les sémaphores sont initialisé à 0 car par défaut il n'y a rien à lire
     * et on n'a rien à écrire
     */
    readAvail = new Semaphore("read avail", 0);
    writeDone = new Semaphore("write done", 0);
    console = new Console (readFile, writeFile, ReadAvail, WriteDone, 0);
}

SynchConsole::~SynchConsole() {
    delete console;
    delete writeDone;
    delete readAvail;
    delete mutex;
}

void SynchConsole::SynchPutChar(const char ch) {
    /* On écrit un char et on se bloque en attendant que le traitant appelle
     * (WriteDone>V())
     */
    console->PutChar(ch);
    writeDone->P();
}

char SynchConsole::SynchGetChar() {
    /* Lorsqu'il y a rien à lire, on se bloque, et dès qu'il y a quelque chose
     * à lire, on sait qu'on sera débloqué (ReadAvail->V())
     */
    readAvail->P();
    return console->GetChar();
}

void SynchConsole::SynchPutString(const char string[]) {
}

void SynchConsole::SynchGetString(char *buffer, int n) {
}
\end{lstlisting}

On peut tester la synchconsole depuis le fichier
\textit{code/userprog/synchconsole.h}

\begin{lstlisting}
void SynchConsoleTest(char *in, char*out)
{
    char ch;
    synchconsole = new SynchConsole(in, out);
    while ((ch = synchconsole->SynchGetChar()) != EOF )
      synchconsole->SynchPutChar(ch);
    fprintf(stderr, "Solaris: EOF detected in SynchConsole!\n");
}
\end{lstlisting}

\section{Appels systèmes}
Dans cette partie on va implémenter les appels systeme qui utiliseront notre
synconsole.
\subsection{PutChar}
\subsubsection{Mise en place}
On commence par déclarrer la fonction d'appel système \textit{void PutChar(char
c)} et son numéro

\textit{code/userprog/syscall.h}
\begin{lstlisting}
#define SC_PutChar  11
// ...
void PutChar(char c);
\end{lstlisting}

Un appel système entraine un changement d'environement, du mode au mode noyau.
Il faut donc écrire (en assembleur) le code qui permet de faire cette
interuption pour basculer en mode noyau et prévoir le retour au programme
utilisateur :

\textit{code/test/start.S}
\begin{lstlisting}
    .globl PutChar
    .ent PutChar
PutChar:
    /* On place le signal dans le registre R2
        Il va servir au handler d'exceptions pour qu'il puisse
        savoir qui sera le traitant de cette exception.
    */
    addiu $2,$0,SC_PutChar
     /* syscall provoque un déroutement et place le compteur de
        programme (PC) à la première instruction du traitant :
        ExceptionHandler
        */
    syscall
     /* Maintenant on revient au programme appellant
        Le registre R31 sauvegarde l'adresse de retour de la
        fonction appellante
        */
    j    $31
    .end PutChar
\end{lstlisting}

Le traitrement sera fait dans \textit{code/userprog/exception.cc}. Pour
faciliter l'ajout de nouveaux appels systèmes on utilisera un switch/case qui
associe un traitrement a chaque numéro d'appel système.

\textit{code/userprog/exception.cc}
\begin{lstlisting}
void ExceptionHandler (ExceptionType which)
{
    int type = machine->ReadRegister (2);

    if (which == SyscallException) {
        switch (type) {

        case SC_Halt: {
          DEBUG('a', "Shutdown, initiated by user program.\n");
          interrupt->Halt();
          break;
        }

        case SC_PutChar: {
          DEBUG('a', "PutChar, initiated by user program.\n");
          // On récupère le premier parametre
          char c = (char)(machine->ReadRegister(4));
          // on l'affiche dans grâce a synchonsole
          synchconsole->SynchPutChar(c);
          break;
        }

        default: {
          printf("Unexpected user mode exception %d %d\n", which, type);
          ASSERT(FALSE);
        }
      }
    }

    // LB: Do not forget to increment the pc before returning!
    UpdatePC ();
    // End of addition
}
\end{lstlisting}

L'objet synchconsole appartient au noyau, il faut donc l'initialiser dans le
fichier \textit{code/threads/system.cc} comme suit

\begin{lstlisting}
#ifdef USER_PROGRAM        // requires either FILESYS or FILESYS_STUB
Machine *machine;        // user program memory and registers
SynchConsole *synchconsole; // On déclare la synchconsole
#endif
// [..]
void Initialize (int argc, char **argv)
{
// ...
#ifdef USER_PROGRAM
machine = new Machine (debugUserProg);    // this must come first
synchconsole = new SynchConsole(NULL, NULL); // On déclare la synchconsole sans
aucun in/out
#endif
// ..
}

void Cleanup ()
{
// ..
#ifdef USER_PROGRAM
delete machine;
delete synchconsole;
#endif
// ..
}
\end{lstlisting}

\textit{code/threads/system.h}
\begin{lstlisting}
#ifdef USER_PROGRAM
#include "machine.h"
#include "../userprog/synchconsole.h"
extern Machine *machine;    // user program memory and registers
extern SynchConsole *synchconsole;
#endif
\end{lstlisting}

Maintenant on peut utiliser synchconsole en important \textit{system.h}.

\subsubsection{Test de PutChar()}

\textit{code/userprog/putchar.c}
\begin{lstlisting}
#include "syscall.h"

int main() {
    PutChar('a');
    PutChar('\n');
    Halt();
}
\end{lstlisting}

\begin{lstlisting}
$ ./build-origin/nachos-userprog -x ./build/putchar
a
Machine halting!

Ticks: total 253, idle 200, system 30, user 23
Disk I/O: reads 0, writes 0
Console I/O: reads 0, writes 2
Paging: faults 0
Network I/O: packets received 0, sent 0

Cleaning up...
\end{lstlisting}

\subsubsection{Terminaison}
Un programme est obligé d'appeller Halt() pour dire qu'il s'est terminé, ce qui
n'est pas pratique en temps normal. Un programme est lancé par la méthode
\textit{Machine::Run()} qui ne termine pas :

\textit{code/machine/mipssim.cc}
\begin{lstlisting}
void Machine::Run()
{
    // [...]
    interrupt->setStatus(UserMode);
    for (;;) {
        OneInstruction(instr);
        interrupt->OneTick();
        if (singleStep && (runUntilTime <= stats->totalTicks))
            Debugger();
    }
}
\end{lstlisting}

L'absence de \textit{Halt()} provoque une interuption :

\begin{lstlisting}
./build-origin/nachos-userprog -x ./build/putchar
a
Unexpected user mode exception 1 1
\end{lstlisting}

L'interuption 1 correspond à l'appel système \textit{Exit()}, il faut
donc implementer cet appel système, qui se contentera, dans un premier temps,
d'eteindre explicitement la machine:

\textit{code/userprog/exception.cc}
\begin{lstlisting}
        case SC_Exit: {
          DEBUG('a', "Exit, initiated by user program.\n");
          interrupt->Halt();
          break;
        }
\end{lstlisting}

\begin{lstlisting}
$ ./build-origin/nachos-userprog -x ./build/putchar
a
Machine halting!

Ticks: total 260, idle 200, system 30, user 30
Disk I/O: reads 0, writes 0
Console I/O: reads 0, writes 2
Paging: faults 0
Network I/O: packets received 0, sent 0

Cleaning up...
\end{lstlisting}

\subsection{PutString}
La difference entre PutChar et PutString c'est que PutString prend une pointeur
sur une chaine de caractères en mémoire \textbf{user}. On va donc devoir par
précaution, préalablement la copier dans un buffer en mémoire \textbf{noyau}.

\textit{code/userprog/exception.cc}
\begin{lstlisting}
char * ReadStringFromMachine(int from, unsigned max_size) {
  /* On copie octet par octet, de la mémoire user vers la mémoire noyau (buffer)
   * en faisant attention à bien convertir explicitement en char
   */
  int byte;
  unsigned int i;
  char * buffer = new char[max_size];
  for(i = 0; i < max_size-1; i++) {
    machine->ReadMem(from+i,1, &byte);
    if((char)byte=='\0')
      break;
    buffer[i] = (char) byte;
  }
  buffer[i] = '\0';
  return buffer;
}
\end{lstlisting}

La mie en place de l'appel systeme est similaire à \textit{PutChar()}. Il faut
simplement spécifier la taille du buffer de copie.

\textit{code/userprog/exception.cc}
\begin{lstlisting}
        case SC_PutString: {
          DEBUG('a', "PutString, initiated by user program.\n");
          // Le premier argument (registre R4) c'est l'adresse de la chaine de
caractere
          // Que l'ont recopie dans le monde linux (noyau)
          // R4 >> pointeur vers la mémoire  MIPS
          // MAX_STRING_SIZE est defni prealablement dans code/threads/system.h
          char *buffer = ReadStringFromMachine(machine->ReadRegister(4),
MAX_STRING_SIZE);
          synchconsole->SynchPutString(buffer);
          delete [] buffer;
          break;
        }
\end{lstlisting}


Dans la SynchConsole on implemente \textit{SynchConsole::SynchPutString} :

\textit{code/userprog/synchonsole.cc}
\begin{lstlisting}
void SynchConsole::SynchPutString(const char string[]) {
    /* On utilise un mutex pour que les appels SynchPutString soient atomiques
     * C'est à dire que deux appels à SynchPutString() affichent correctement
     * les chaines de caractères...
     * * */
    mutex->P();
    for(int i=0; i<MAX_STRING_SIZE-1;i++) {
        if(string[i] == '\0')
            break;
        this->SynchPutChar(string[i]);
    }
    mutex->V();
}
\end{lstlisting}

On utilise un mutex (Sémaphore initialisé à 1) pour assuré l'atomicité de
PutString. En effet, on souhaite avoir tous les caractères dans le bon ordre,
et deux appels a PutString doivent se faire l'un après l'autre.

\subsection{GetChar GetString}

Ces deux appels systèmes sont symétriques à PutChar et PutString. Dans le cas
de GetChar, rien de plus simple, etant donné qu'il renvoi directement la
valeur (C'est le registre 2 qui est utilisé pour les valeurs de retour). On
ecrit directement dans le registre la valeur de retour de SynchGetChar().

\textit{code/userprog/exception.cc}
\begin{lstlisting}
        case SC_GetChar: {
          DEBUG('a', "GetChar, initiated by user program.\n");
          machine->WriteRegister(2,(int) synchconsole->SynchGetChar());
          break;
        }
\end{lstlisting}

Pour GetString, on ecrit dans un buffer intermediaire, puis on copie ce buffer
dans la mémoire user à l'adresse donnée a l'appel système.

\begin{lstlisting}
        case SC_GetString: {
          DEBUG('a', "GetString, initiated by user program.\n");

          // le premier argument est une adresse (char *)
          int to = machine->ReadRegister(4);
          // le second est un int >> la taille
          int size = machine->ReadRegister(5);
          // On donne pas acces à la mémoire directement, on ecrit ecrit dans un
buffer
          // Peut etre pas obligé, mais au cas ou on utilise un buffer..
          char * buffer = new char[MAX_STRING_SIZE];
          synchconsole->SynchGetString(buffer, size);
          WriteStringToMachine(buffer, to, size);
          delete [] buffer;
          break;
        }
\end{lstlisting}


\begin{lstlisting}
void WriteStringToMachine(char * string, int to, unsigned max_size) {
  /* On copie octet par octet, en faisant attention à bien convertir
   * explicitement en char
   */
  char * bytes = (char *)(&machine->mainMemory[to]);
  for(unsigned int i = 0; i < max_size-1; i++) {
    bytes[i] = string[i];
    if(string[i]=='\0')
      break;
  }
}
\end{lstlisting}

\textit{code/userprog/synchonsole.cc}
\begin{lstlisting}
void SynchConsole::SynchGetString(char *buffer, int n) {
  /* On utilise un mutex pour que tous les appels SynchGetString soient
   * atomiques.
   * * */
  int i;
  char c;
  mutex->P();
  for (i=0; i<n-1; i++) {
    c = this->SynchGetChar();
    // CTRL+D pour arrêter la saisie
    if(c == EOF)
      break;
    else
      buffer[i] = c;
  }
  buffer[i] = '\0';
  mutex->V();
}
\end{lstlisting}

\subsection{PutInt et GetInt}

Pour PutInt et GetInt on va utiliser les fonction \textbf{sscanf} et
\textbf{snprintf}. Pour facilier la saisie, on ajoute la fonction
\textit{SynchConsole::SynchGetString(char *buffer, int n, char delim)} qui
permet de lire un chaine de caractere et de s'arreter des qu'on rencontre un
délimiteur (delim). Dans notre cas, on va utiliser '\textbackslash{n}'
comme délémiteur lors de la saisie de nombre entier :

\textit{code/userprog/synchonsole.cc}
\begin{lstlisting}
void SynchConsole::SynchPutInt(int value) {
  char * buffer = new char[MAX_STRING_SIZE];
  // on ecrit dans le buffer la valeur avec sprintf
  snprintf(buffer,MAX_STRING_SIZE, "%d", value);
  this->SynchPutString(buffer);
  delete [] buffer;
}

int SynchConsole::SynchGetInt() {
  int value;
  char * buffer = new char[MAX_STRING_SIZE];
  this->SynchGetString(buffer, MAX_STRING_SIZE, '\n');
  sscanf(buffer, "%d", &value);
  delete [] buffer;
  return value;
}
\end{lstlisting}

\textit{code/userprog/exception.cc}
\begin{lstlisting}
        case SC_PutInt: {
          DEBUG('a', "PutInt, initiated by user program.\n");
          // le premier est la valeur int
          int value = machine->ReadRegister(4);
          synchconsole->SynchPutInt(value);
          break;
        }

        case SC_GetInt: {
          DEBUG('a', "GetInt, initiated by user program.\n");
          int value = synchconsole->SynchGetInt();
          machine->WriteRegister(2, value);
          break;
        }
\end{lstlisting}

\section{Détection de fin de fichier}

flemme...

\section{Test Nachos étape 2}

Voici le programme de test :

\begin{lstlisting}
#include "syscall.h"

int main() {
  PutString("Veuillez saisir un nombre : \n");
  int nombre = GetInt();
  PutString("Nombre +10 = ");PutInt(nombre+10);PutChar('\n');
  PutString("Veuillez saisir une lettre : \n");
  char c = GetChar();
  PutString("Voici la lettre : ");
  PutChar(c);
  PutString("\nVeuillez saisir une phrase (max = 100) : ");
  char buffer[100];
  GetString(buffer,100);
  PutString("\nVoici la phrase : ");
  PutString(buffer);
  PutChar('\n');
  return 0;
}
\end{lstlisting}

\begin{lstlisting}
$ ./build-origin/nachos-userprog -x ./build/etape2
Veuillez saisir un nombre :
5
Nombre +10 = 15
Veuillez saisir une lettre :
g
Voici la lettre : g
Veuillez saisir une phrase (max = 100) : le mot de la fin

Voici la phrase :
le mot de la fin

Machine halting!

Ticks: total 1136142117, idle 1136139819, system 2170, user 128
Disk I/O: reads 0, writes 0
Console I/O: reads 22, writes 174
Paging: faults 0
Network I/O: packets received 0, sent 0

Cleaning up...
\end{lstlisting}


\end{document}

